% !PRINTED 97
\weberpage{97}

Die Transformationen $S$ setzen sich also nach derselben Regel zusammen wie die linearen Substitutionen und Matrizes, die wir im sechsten Abschnitte des zweiten Bandes betrachtet haben. Der Grad einer zusammengesetzten Transformation ist gleich dem Produkte der Grade der Komponenten. Diese Matrizes sind hier an die Voraussetzung gebunden, daß ihre Elemente ganze Zahlen und ihre Determinante positiv ist.

Diese Eigenschaften bleiben bei der Zusammensetzung der Transformationen erhalten. Trotzdem bildet die Gesamtheit $\mathfrak{S}$ der Transformationen $S$ keine Gruppe, so wenig wie die Gesamtheit der natürlichen Zahlen bei der Komposition durch Multiplikation eine Gruppe ist; denn es läßt sich bei gegebenem $S'$, $S''$ nicht immer ein $S$ bestimmen, das der Bedingung (6) genügt, was doch (nach Bd. II, \S\ 1, 4.) für eine Gruppe erforderlich wäre.

Durch die spezielle Transformation vom Grade $m^2$:
\[
M = \begin{pmatrix} \pm m, & 0 \\ 0, & \pm m \end{pmatrix}
\]
gehen die Perioden $\omega_1, \omega_2$ in $\omega_1' = \omega_1/m$, $\omega_2' = \omega_2/m$ über, und das Periodenverhältnis $\omega = \omega_2/\omega_1$ bleibt ungeändert. Diese Transformationen heißen \gesperrt{Multiplikationen} (Ähnlichkeits-Transformationen, Bd. II, \S\ 41). Es ist darunter die identische Substitution
\[
1 = \begin{pmatrix} 1, & 0 \\ 0, & 1 \end{pmatrix}
\]
enthalten, die alles ungeändert läßt und bei der Komposition die Rolle der Einheit spielt.

\gesperrt{Die Multiplikationen sind bei der Zusammensetzung mit jeder Transformation $S$ vertauschbar:}

\weq{7}{SM = MS.}

Hält man in $SM$ oder $MS$ die Transformation $S$ fest und läßt $M$ die Gesamtheit $\mathfrak{M}$ der Multiplikationen durchlaufen, so erhält man ein System $\mathfrak{M}S$, das man nach Bd. II, \S\ 46 als eine Kollineation zu bezeichnen hätte. Gehören $S_1$ und $S_2$ einer Kollineation $C$ an und $S_1'$ und $S_2'$ einer Kollineation $C'$, so gehören auch $S_1'S_1$ und $S_2'S_2$ derselben Kollineation $C''$ an. Man kann so, indem man $C'' = C'C$ setzt, die Kollineationen zusammensetzen. Bei dieser Zusammensetzung spielt die Kolli-
%CONT>>
